\documentclass{zhvt-classic}

% 使用本地源流明体（GenRyuMin2 TW）按文件路径加载，不依赖系统字体名
% 版框（内框／界栏／书口侧栏／鱼尾）改用与版心印章同一印泥棕色：
%   印章图片 yshs.png 实色 = RGB(135,68,52) = HTML 874434
\usepackage{xcolor}
\usepackage{tikz}
% sealbrown 双注册：xcolor 供 TikZ/pgf，L3 供 cls 版框
\definecolor{sealbrown}{HTML}{874434}
\ExplSyntaxOn
\color_set:nnn { sealbrown } { HTML } { 874434 }
\ExplSyntaxOff

\zhvtset{
  grid color       = sealbrown,
  main font        = GenRyuMin2TW-SB,
  main font path   = C:/Users/Administrator/texmf/fonts/opentype/public/GenRyuMin2TW/,
jiazhu font      = GenRyuMin2TW-M,
  jiazhu font path = C:/Users/Administrator/texmf/fonts/opentype/public/GenRyuMin2TW/,
  % 页码框容纳 4 字：正常页「九十七」+前置空 1 格=4 行恰好；
  % 100+ 页「一百零一」共 5 行超出，末字「一/二/三」被框裁掉（2026-10-04 实测），故放宽到 5
  shukou page number chars = 5,
  % 书口书名框高从默认 6 字减到 4 字，把省出的 2 字高拨给 chapter 框，
  % 使 6 字书名+卷名放进 9 字高的 chapter 框后不再压线（2026-10-05）
  shukou title chars = 4,
}

\title{史记}[田儋列传]
\author{汉·司马迁}
\maker{闲叟漫记·重製}


% ===== 索隐述赞样式（2026-11-08 四次定稿 · 愚大人亲手改定参数）=====
% ⚠ 竖排前提：自定义字体家族必须带 RawFeature={vertical:+vert}；
%   simhei.ttf 加 +vhal 会整体加载失败静默回退，只写 +vert；simfang 同理。
\newCJKfontfamily{\shuzanfont}[Path=C:/Windows/Fonts/,RawFeature={vertical:+vert}]{simhei.ttf}
\newCJKfontfamily{\shuzanbody}[Path=C:/Windows/Fonts/,RawFeature={vertical:+vert}]{simfang.ttf}
% 题头：152x30pt 圆角 6pt 棕底白字；baseline (0,2) 使整框落入界行格子不压版框线；
% 白字节点 x 不补偿、y +8.2pt（像素级墨迹实测 dx/dy≈0，L/R 留白约 20pt）。
\newcommand{\shuzanhead}[1]{%
  \begin{tikzpicture}[baseline={(0,2)}]
    \fill[sealbrown, rounded corners=6pt]
      (-76pt,-15pt) rectangle (76pt,15pt);
    \node[text=white, inner sep=0pt, anchor=center] at (0pt,8.2pt)
      {\shuzanfont\fontsize{23}{23}\selectfont #1};
  \end{tikzpicture}%
}
% ★ 修正书口卷名位置：chapter 框起始位置改回 cls 默认值，避免全局影响。
%   只在每个卷首页通过 \juan 命令里的空命令微调该页位置（安全，不会跨页污染）。
\ExplSyntaxOn
\AtBeginDocument{
  \dim_set:Nn \l__zhvt_shukou_chapter_title_start_dim
    { \l__zhvt_shukou_title_height_dim + \l__zhvt_font_size_dim * 2 }
  \dim_set:Nn \l__zhvt_shukou_page_number_start_dim
    { \l__zhvt_shukou_chapter_title_start_dim + \l__zhvt_shukou_chapter_title_height_dim }
}
\ExplSyntaxOff

\begin{document}

% ★★ 版心（书口）排布覆写：把书名放到鱼尾之下（2026-10-04）
%
%   cls 默认版心自上而下：书名（title 静态框，在鱼尾之上）→ 鱼尾 → 篇名（chapter 动态框）→ 页码
%   本模板改为：          鱼尾（位置不动）→ 书名 → 篇名 → 页码分割线 → 页码
%
%   手法：把 cls 放在 title 静态框里的书名，改放进 chapter 动态框（鱼尾之下）；
%         title 静态框留空。书名取 \l__zhvt_title_str，即 \title{...} 的第一参数，
%         所以换书时只改 \title，不必在此重复写书名。
%
%   ⚠ 三条必须遵守的写法（2026-10-04 实测）：
%     ① 汉字必须用 \hbox{...} 包住。flowfram 会把框内容装进一个宽 14.4mm（= 一个 line_sep
%        = 书口栏宽）的纵向 minipage 并按该宽度折行，而整页靠 PDF /Rotate 90 显示、
%        框内文字本身不变向 —— 不包 \hbox 时「围炉夜话」4 字会被逐字折成 4 行，
%        在版面上就成了横排 4 字并横向溢出版框。
%     ② 竖排下**一个 \hbox 就是一整列**。书名与篇名要同处一列，必须放进同一个 \hbox；
%        分成两个 \hbox 时，后一个会被当成新的一列、甩到隔壁栏去
%        （实测「序」跑到书口栏外侧 14.4mm 的正文列上）。中间的 \quad 在框内是横向空白，
%        显示出来正好是书名与篇名之间的纵向间隔。
%     ③ ★ 内容盒要写成 **与框等宽的定宽盒** `\hbox to \l__zhvt_shukou_chapter_title_height_dim{…\hfil}`。
%        实测：chapter 框是把内容盒沿纵向**居中**摆放的，内容宽度一变，书名就跟着漂：
%        自然宽度的 \hbox（书名 38.4mm）落在 x=157.2mm，加篇名后（57.6mm）落在 x=152.4mm，
%        目录页（无篇名）与正文页的书名就差半格。做成等宽盒后，无论有没有篇名，
%        内容固定从框顶起排，实测两页均落在 x=142.8mm，不再漂。
\ExplSyntaxOn
% 把书口 chapter 框（书名+篇名）起始位置下移 1 字高，避免 4 字卷名顶到鱼尾
% ⚠ cls 的 \__zhvt_calculate_dimensions: 在 \AtBeginDocument 里执行，会重算该 dim；
%   必须把覆写挂到 begin document 之后，否则会被二次覆盖（2026-10-05 实测）。
% (旧方案已移到导言区，直接修改 chapter 框起始位置)
% 卷标题命令 #1=书口篇名(鱼尾) #2=正文卷名（含卷序，卷名内空格用 ~）
%
% ★★ 目录页两项修订（2026-10-05，愚大人指出）★★
%   ① 页码起点统一对齐：
%      原状 —— 目录项写成 \par\hspace{4em}卷名\hspace{4em}页码，页码起点 = 卷名实宽 + 4em。
%      卷名实宽随字数变：「卷一~梁惠王上」6 字 +~、「卷七~离娄上」5 字 +~，
%      差 1 个字位 → 卷七~卷十 4 条的页码整体比其余 10 条高 27.21pt（一个字），看着上下不齐。
%      改法 —— 卷名补全角空格（宽 = 1 字 = 27.215pt）补足到 **7 字符位**再写进目录，
%      14 条页码起点遂全部落在同一条竖线上（起点对齐；末字随页码字数自然参差）。
%      ⚠ 只对目录里的副本 \l_tmpc_tl 补，正文卷名仍用 \l_tmpa_tl 原样 —— 正文像素不受影响。
%   ② 卷名整体上移半字位：前导 \hspace 由 4em 改 3.5em。
%      实测原状目录块在版框内偏下（天头留白 139.6pt / 地脚留白 110.7pt），
%      上移 0.5em = 13.6pt 后上下基本均分，视觉居中。
\NewDocumentCommand { \juan } { mm } {
	\clearpage
	\tl_set:Nn \l__zhvt_chapter_title_tl { #1 }
	\tl_set:Nn \l_tmpa_tl { #2 }
	\tl_set_eq:NN \l_tmpc_tl \l_tmpa_tl
	\int_step_inline:nn { 7 - \str_count:n { #2 } }
	  { \tl_put_right:Nn \l_tmpc_tl { 　 } }
	\tl_set:Nx \l_tmpb_tl { \zhnum { page } }
	\tl_put_left:Nn \l_tmpb_tl { \hspace { 4em } }
	\tl_put_left:Nn \l_tmpb_tl { \l_tmpc_tl }
	\tl_put_left:Nn \l_tmpb_tl { \par \hspace { 3.5em } }
	\addtocontents { toc } { \l_tmpb_tl }
	\skip_horizontal:n { \l__zhvt_font_size_dim * 4 }
	\l_tmpa_tl
	\par
	\thispagestyle { empty }
}
\cs_set_protected:Npn \__zhvt_set_shukou_content: {
	% ★ 鱼尾上方留白（title 静态框）嵌入闲叟钤印：
	%   整页 /Rotate 90，故此处把图用 \rotatebox{90} 转正为竖章；
	%   尺寸卡在版心条幅内（书内约 7mm 宽 × 33.8mm 高，留余量不压鱼尾）。
	\setstaticcontents * { title } {
		\vbox_to_ht:nn { \l__zhvt_baseline_skip_dim } {
			\vskip 14.2pt
			\hbox_to_wd:nn { \l__zhvt_shukou_title_height_dim } {
				\hfil
				\rotatebox [ origin = c ] { 90 } { \includegraphics [ width = 7mm , height = 33.8mm ] { yshs.png } }
				\hfil
			}
		}
	}
% ★ 版心（书口）统一改法（2026-11-16 愚大人定）：18pt＋框内居中＋书口文字同版框色
%   ① 原 27.2pt 时 7 字长篇名（如「廉颇蔺相如列传」）末字溢出 chapter 框、压住页码；
%      改 \fontsize{18pt}{18pt}，前后 \hfil 令字串在框内居中（不随篇名字数漂移）。
%   ② \vbox{\vskip 8.2pt} 定量下移，抵消缩字号后栏宽方向 8.2pt 上偏（像素实测 y 中心 543.3 = 原版）。
%      ⚠ \raisebox / \vspace* 在 flowfram 框内无效；换字号须同步换补偿：16pt→10.1、19pt→7.4。
%   ③ chapter 与 page~number 均置 \color{sealbrown}（HTML 874434），与版框／界行同色；正文仍黑。
%   ⚠ 页码字位、钤印、鱼尾、界行、正文一律未动；五部合订本同此口径。
	\setdynamiccontents * { chapter }
	{ \color{sealbrown} \vbox { \vskip 8.2pt \hbox to \l__zhvt_shukou_chapter_title_height_dim
		{ \hfil \fontsize{18pt}{18pt}\selectfont \l__zhvt_title_str \quad \l__zhvt_chapter_title_tl \hfil } } }
	\setdynamiccontents * { page~number } { \color{sealbrown} \quad \zhnum { page } }
}
\__zhvt_set_shukou_content:
\ExplSyntaxOff

\maketitle

% 单篇，无多卷目录，不渲染目录页

\juan{田儋列传}{卷九十四~田儋列传第三十四}
\hspace*{2em}田儋者。狄人也。故齐王田氏族也。儋从弟田荣。荣弟田横。皆豪。宗彊。能得人。\\
\hspace*{2em}陈涉之初起王楚也。使周市略定魏地。北至狄。狄城守。田儋详为缚其奴。从少年之廷。欲谒杀奴。见狄令。因击杀令。而召豪吏子弟曰。﹁诸侯皆反秦自立。齐。古之建国。儋。田氏。当王。﹂遂自立为齐王。发兵以击周市。周市军还去。田儋因率兵东略定齐地。\\
\hspace*{2em}秦将章邯围魏王咎於临济。急。魏王请救於齐。齐王田儋将兵救魏。章邯夜衔枚击。大破齐。魏军。杀田儋於临济下。儋弟田荣收儋馀兵东走东阿。\\
\hspace*{2em}齐人闻王田儋死。乃立故齐王建之弟田假为齐王。田角为相。田间为将。以距诸侯。\\
\hspace*{2em}田荣之走东阿。章邯追围之。项梁闻田荣之急。乃引兵击破章邯军东阿下。章邯走而西。项梁因追之。而田荣怒齐之立假。乃引兵归。击逐齐王假。假亡走楚。齐相角亡走赵。角弟田间前求救赵。因留不敢归。田荣乃立田儋子市为齐王。荣相之。田横为将。平齐地。\\
\hspace*{2em}项梁既追章邯。章邯兵益盛。项梁使使告赵。齐。发兵共击章邯。田荣曰。﹁使楚杀田假。赵杀田角。田间。乃肯出兵。﹂楚怀王曰。﹁田假与国之王。穷而归我。杀之不义。﹂赵亦不杀田角。田间以市於齐。齐曰。﹁蝮螫手则斩手。螫足则斩足。何者。为害於身也。今田假。田角。田间於楚。赵。非直手足戚也。何故不杀。且秦复得志於天下。则齮龁用事者坟墓矣。﹂楚。赵不听。齐亦怒。终不肯出兵。章邯果败杀项梁。破楚兵。楚兵东走。而章邯渡河围赵於钜鹿。项羽往救赵。由此怨田荣。\\
\hspace*{2em}项羽既存赵。降章邯等。西屠咸阳。灭秦而立侯王也。乃徙齐王田市更王胶东。治即墨。齐将田都从共救赵。因入关。故立都为齐王。治临淄。故齐王建孙田安。项羽方渡河救赵。田安下济北数城。引兵降项羽。项羽立田安为济北王。治博阳。田荣以负项梁不肯出兵助楚。赵攻秦。故不得王。赵将陈馀亦失职。不得王。二人俱怨项王。\\
\hspace*{2em}项王既归。诸侯各就国。田荣使人将兵助陈馀。令反赵地。而荣亦发兵以距击田都。田都亡走楚。田荣留齐王市。无令之胶东。市之左右曰。﹁项王彊暴。而王当之胶东。不就国。必危。﹂市惧。乃亡就国。田荣怒。追击杀齐王市於即墨。还攻杀济北王安。於是田荣乃自立为齐王。尽并三齐之地。\\
\hspace*{2em}项王闻之。大怒。乃北伐齐。齐王田荣兵败。走平原。平原人杀荣。项王遂烧夷齐城郭。所过者尽屠之。齐人相聚畔之。荣弟横。收齐散兵。得数万人。反击项羽於城阳。而汉王率诸侯败楚。入彭城。项羽闻之。乃醳齐而归。击汉於彭城。因连与汉战。相距荥阳。以故田横复得收齐城邑。立田荣子广为齐王。而横相之。专国政。政无巨细皆断於相。\\
\hspace*{2em}横定齐三年。汉王使郦生往说下齐王广及其相国横。横以为然。解其历下军。汉将韩信引兵且东击齐。齐初使华无伤。田解军於历下以距汉。汉使至。乃罢守战备。纵酒。且遣使与汉平。汉将韩信已平赵。燕。用蒯通计。度平原。袭破齐历下军。因入临淄。齐王广。相横怒。以郦生卖己。而亨郦生。齐王广东走高密。相横走博。守相田光走城阳。将军田既军於胶东。楚使龙且救齐。齐王与合军高密。汉将韩信与曹参破杀龙且。虏齐王广。汉将灌婴追得齐守相田光。至博。而横闻齐王死。自立为齐王。还击婴。婴败横之军於嬴下。田横亡走梁。归彭越。彭越是时居梁地。中立。且为汉。且为楚。韩信已杀龙且。因令曹参进兵破杀田既於胶东。使灌婴破杀齐将田吸於千乘。韩信遂平齐。乞自立为齐假王。汉因而立之。\\
\hspace*{2em}後岁馀。汉灭项籍。汉王立为皇帝。以彭越为梁王。田横惧诛。而与其徒属五百馀人入海。居岛中。高帝闻之。以为田横兄弟本定齐。齐人贤者多附焉。今在海中不收。後恐为乱。乃使使赦田横罪而召之。田横因谢曰。﹁臣亨陛下之使郦生。今闻其弟郦商为汉将而贤。臣恐惧。不敢奉诏。请为庶人。守海岛中。﹂使还报。高皇帝乃诏卫尉郦商曰。﹁齐王田横即至。人马从者敢动摇者致族夷。﹂乃复使使持节具告以诏商状。曰。﹁田横来。大者王。小者乃侯耳。不来。且举兵加诛焉。﹂田横乃与其客二人乘传诣雒阳。\\
\hspace*{2em}未至三十里。至尸乡厩置。横谢使者曰。﹁人臣见天子当洗沐。﹂止留。谓其客曰。﹁横始与汉王俱南面称孤。今汉王为天子。而横乃为亡虏而北面事之。其耻固已甚矣。且吾亨人之兄。与其弟并肩而事其主。纵彼畏天子之诏。不敢动我。我独不愧於心乎。且陛下所以欲见我者。不过欲一见吾面貌耳。今陛下在洛阳。今斩吾头。驰三十里间。形容尚未能败。犹可观也。﹂遂自刭。令客奉其头。从使者驰奏之高帝。高帝曰。﹁嗟乎。有以也夫。起自布衣。兄弟三人更王。岂不贤乎哉。﹂为之流涕。而拜其二客为都尉。发卒二千人。以王者礼葬田横。\\
\hspace*{2em}既葬。二客穿其冢旁孔。皆自刭。下从之。高帝闻之。乃大惊。大田横之客皆贤。吾闻其馀尚五百人在海中。使使召之。至则闻田横死。亦皆自杀。於是乃知田横兄弟能得士也。\\
\hspace*{2em}太史公曰。甚矣蒯通之谋。乱齐骄淮阴。其卒亡此两人。蒯通者。善为长短说。论战国之权变。为八十一首。通善齐人安期生。安期生尝干项羽。项羽不能用其筴。已而项羽欲封此两人。两人终不肯受。亡去。田横之高节。宾客慕义而从横死。岂非至贤。余因而列焉。不无善画者。莫能图。何哉。\\

% ===== 索隐述赞（司马贞《史记索隐》）· 统一样式 =====
% 题头「索隐述赞曰」五字并入反白框；正文每四句一列，定宽盒 \hfil 居列中；
% \fontsize{24}{41} 列距对齐全书界行（121+41k）；\clearpage 使述赞独占一页。
\clearpage
\shuzanhead{索隐述赞曰}\linebreak
{\shuzanbody\fontsize{24}{41}\selectfont
\hbox to \hsize{\hfil 秦项之际。天下交兵。六国树党。自置豪英。\hfil}\linebreak
\hbox to \hsize{\hfil 田儋殒寇。立市相荣。楚封王假。齐破郦生。\hfil}\linebreak
\hbox to \hsize{\hfil 兄弟更王。海岛传声。\hfil}%
\par}
\end{document}
